\chapter{Documentation des services SOAP}\label{chap:soap}

\section{Principes communs}

Les trois services de référentiel suivent le même modèle.

\begin{itemize}
  \item \textbf{Contrat d'abord.} Chaque WSDL (style \emph{document/literal}, SOAP~1.1) est écrit à la main et documenté opération par opération (\code{wsdl:documentation}) ; il fait foi. Les classes Java sont générées à la compilation par \code{wsimport}, et l'implémentation renvoie au fichier par l'attribut \code{wsdlLocation} de \code{@WebService} : le WSDL publié par le service est exactement celui du dépôt.
  \item \textbf{Espaces de noms.} \code{http://stages.inpt.ma/etudiants}, \code{…/entreprises}, \code{…/offres} ; l'en-tête de corrélation est dans \code{http://stages.inpt.ma/commun}.
  \item \textbf{Adresses.} \code{http://localhost:8081/ws/etudiants}, \code{:8082/ws/entreprises}, \code{:8083/ws/offres} ; le WSDL s'obtient en ajoutant \code{?wsdl} (lecture sans authentification).
  \item \textbf{Validation.} Les données reçues sont validées (Bean Validation : formats du matricule, de l'ICE, du téléphone, longueur des champs, valeurs énumérées) avant tout accès à la base ; chaque champ invalide produit un message en français dans le détail de la faute.
  \item \textbf{Fautes typées.} Toute erreur prévue est une faute déclarée dans le WSDL, avec un détail structuré : un \code{code} stable, exploitable par un programme, un \code{message} en français et une liste de \code{details}. Aucune trace d'exécution Java n'est renvoyée.
  \item \textbf{Sécurité.} Chaque appel doit porter un en-tête WS-Security \emph{UsernameToken} (section~\ref{sec:securite-soap}).
\end{itemize}

\begin{table}[H]
\centering\small
\begin{tabularx}{\linewidth}{@{}l>{\raggedright\arraybackslash}X l@{}}
\toprule
\textbf{Faute} & \textbf{Situation} & \textbf{Exemples de codes} \\ \midrule
\code{*IntrouvableFault} & L'identifiant ne correspond à aucun étudiant, entreprise ou offre. & \code{ETUDIANT\_INTROUVABLE} \\
\code{DonneesInvalidesFault} & Données absentes ou mal formées. & \code{DONNEES\_INVALIDES} \\
\code{ConflitFault} & L'opération contredit l'état ou une règle de gestion. & \code{TRANSITION\_INTERDITE}, \code{PLUS\_DE\_PLACE} \\
\code{ServiceIndisponibleFault} & La base ou un service dépendant ne répond pas. & \code{SERVICE\_OFFRES\_INDISPONIBLE} \\
\code{ErreurInterneFault} & Erreur inattendue ; le détail technique reste dans les journaux. & \code{ERREUR\_INTERNE} \\
\bottomrule
\end{tabularx}
\caption{Fautes SOAP communes aux trois services}\label{tab:fautes}
\end{table}

\section{Opérations}

Les tableaux suivants sont produits automatiquement à partir des WSDL du dépôt (script \code{rapport/outils/generer\_tables.py}) : ils ne peuvent pas diverger du contrat. La dernière colonne indique les fautes métier attendues pour chaque opération ; toutes peuvent en plus renvoyer \code{ServiceIndisponibleFault} et \code{ErreurInterneFault}.

\subsection{S1 Gestion des étudiants}
Le service gère le profil des étudiants et leurs compétences déclarées. \code{verifierEtudiant} est l'opération appelée par le processus à l'étape~2 : elle ne lève pas de faute pour un étudiant inconnu, mais répond \code{existe=false}, ce qui permet au processus de rejeter proprement la candidature (règle R1) au lieu de partir en erreur technique.
{\small% Généré par rapport/outils/generer_tables.py à partir de contracts/src/main/resources/wsdl/etudiants.wsdl. Ne pas modifier.
\begin{longtable}{@{}>{\raggedright\arraybackslash}p{0.27\linewidth}>{\raggedright\arraybackslash}p{0.22\linewidth}>{\raggedright\arraybackslash}p{0.21\linewidth}>{\raggedright\arraybackslash\footnotesize}p{0.22\linewidth}@{}}
\caption{Opérations SOAP de S1 Étudiants (port 8081, 7 opérations)}\label{tab:soap-etudiants}\\
\toprule \textbf{Opération} & \textbf{Entrée} & \textbf{Sortie} & \textbf{Fautes métier} \\ \midrule
\endfirsthead
\toprule \textbf{Opération} & \textbf{Entrée} & \textbf{Sortie} & \textbf{Fautes métier} \\ \midrule
\endhead
\bottomrule
\endlastfoot
\textbf{\texttt{ajouterEtudiant}}\newline {\footnotesize Crée un étudiant.} & \texttt{etudiant} : Etudiant & \texttt{etudiantCree} : Etudiant & \texttt{DonneesInvalides} \newline \texttt{Conflit} \\ \midrule
\textbf{\texttt{consulterEtudiant}}\newline {\footnotesize Renvoie le profil d'un étudiant.} & \texttt{id} : long & \texttt{etudiant} : Etudiant & \texttt{EtudiantIntrouvable} \\ \midrule
\textbf{\texttt{listerEtudiants}}\newline {\footnotesize Liste les étudiants selon des critères facultatifs.} & \texttt{filtre} : FiltreEtudiants, facultatif & \texttt{etudiant} : Etudiant[], facultatif & \texttt{DonneesInvalides} \\ \midrule
\textbf{\texttt{modifierEtudiant}}\newline {\footnotesize Remplace le profil. Un changement de statut suit les mêmes règles que changerStatutEtudiant.} & \texttt{id} : long \newline \texttt{etudiant} : Etudiant & \texttt{etudiantModifie} : Etudiant & \texttt{EtudiantIntrouvable} \newline \texttt{DonneesInvalides} \newline \texttt{Conflit} \\ \midrule
\textbf{\texttt{supprimerEtudiant}}\newline {\footnotesize Supprime un étudiant, sauf pendant un stage.} & \texttt{id} : long & – & \texttt{EtudiantIntrouvable} \newline \texttt{Conflit} \\ \midrule
\textbf{\texttt{verifierEtudiant}}\newline {\footnotesize Indique si l'étudiant existe, son statut et s'il peut candidater (R1). Un étudiant inconnu donne existe=false, pas une faute.} & \texttt{id} : long & \texttt{verification} : VerificationEtudiant & \texttt{ServiceIndisponible} \\ \midrule
\textbf{\texttt{changerStatutEtudiant}}\newline {\footnotesize Change le statut selon les transitions autorisées.} & \texttt{id} : long \newline \texttt{statut} : ACTIF | SUSPENDU | EN\_STAGE | DIPLOME & \texttt{etudiant} : Etudiant & \texttt{EtudiantIntrouvable} \newline \texttt{DonneesInvalides} \newline \texttt{Conflit} \\
\end{longtable}
}

\subsection{S2 Gestion des entreprises}
Le service gère les entreprises partenaires et leur responsable de contact. \code{consulterOffresEntreprise} appelle S3 : S2 ne stocke aucune offre. La suppression d'une entreprise qui a publié des offres est refusée (\code{ConflitFault}) ; il faut la passer en \statut{INACTIVE}.
{\small% Généré par rapport/outils/generer_tables.py à partir de contracts/src/main/resources/wsdl/entreprises.wsdl. Ne pas modifier.
\begin{longtable}{@{}>{\raggedright\arraybackslash}p{0.27\linewidth}>{\raggedright\arraybackslash}p{0.22\linewidth}>{\raggedright\arraybackslash}p{0.21\linewidth}>{\raggedright\arraybackslash\footnotesize}p{0.22\linewidth}@{}}
\caption{Opérations SOAP de S2 Entreprises (port 8082, 7 opérations)}\label{tab:soap-entreprises}\\
\toprule \textbf{Opération} & \textbf{Entrée} & \textbf{Sortie} & \textbf{Fautes métier} \\ \midrule
\endfirsthead
\toprule \textbf{Opération} & \textbf{Entrée} & \textbf{Sortie} & \textbf{Fautes métier} \\ \midrule
\endhead
\bottomrule
\endlastfoot
\textbf{\texttt{ajouterEntreprise}}\newline {\footnotesize Crée une entreprise.} & \texttt{entreprise} : Entreprise & \texttt{entrepriseCreee} : Entreprise & \texttt{DonneesInvalides} \newline \texttt{Conflit} \\ \midrule
\textbf{\texttt{consulterEntreprise}}\newline {\footnotesize Renvoie une entreprise.} & \texttt{id} : long & \texttt{entreprise} : Entreprise & \texttt{EntrepriseIntrouvable} \\ \midrule
\textbf{\texttt{listerEntreprises}}\newline {\footnotesize Liste les entreprises selon des critères facultatifs.} & \texttt{filtre} : FiltreEntreprises, facultatif & \texttt{entreprise} : Entreprise[], facultatif & \texttt{DonneesInvalides} \\ \midrule
\textbf{\texttt{modifierEntreprise}}\newline {\footnotesize Remplace les informations, statut compris.} & \texttt{id} : long \newline \texttt{entreprise} : Entreprise & \texttt{entrepriseModifiee} : Entreprise & \texttt{EntrepriseIntrouvable} \newline \texttt{DonneesInvalides} \newline \texttt{Conflit} \\ \midrule
\textbf{\texttt{supprimerEntreprise}}\newline {\footnotesize Supprime une entreprise sans offre. Interroge S3.} & \texttt{id} : long & – & \texttt{EntrepriseIntrouvable} \newline \texttt{Conflit} \newline \texttt{ServiceIndisponible} \\ \midrule
\textbf{\texttt{verifierEntreprise}}\newline {\footnotesize Indique si l'entreprise est référencée et active (R2). Une entreprise inconnue donne referencee=false, pas une faute.} & \texttt{id} : long & \texttt{verification} : VerificationEntreprise & \texttt{ServiceIndisponible} \\ \midrule
\textbf{\texttt{consulterOffresEntreprise}}\newline {\footnotesize Liste les offres de l'entreprise en appelant S3.} & \texttt{id} : long & \texttt{offre} : OffreResume[], facultatif & \texttt{EntrepriseIntrouvable} \newline \texttt{ServiceIndisponible} \\
\end{longtable}
}

\subsection{S3 Gestion des offres de stage}
Le service gère les offres et porte trois vérifications du processus : la disponibilité (offre active, date limite, places), la compatibilité avec le profil d'un étudiant (filière acceptée, niveau suffisant, score indicatif de compétences) et le retrait atomique d'une place lors d'une acceptation.
{\small% Généré par rapport/outils/generer_tables.py à partir de contracts/src/main/resources/wsdl/offres.wsdl. Ne pas modifier.
\begin{longtable}{@{}>{\raggedright\arraybackslash}p{0.27\linewidth}>{\raggedright\arraybackslash}p{0.22\linewidth}>{\raggedright\arraybackslash}p{0.21\linewidth}>{\raggedright\arraybackslash\footnotesize}p{0.22\linewidth}@{}}
\caption{Opérations SOAP de S3 Offres (port 8083, 8 opérations)}\label{tab:soap-offres}\\
\toprule \textbf{Opération} & \textbf{Entrée} & \textbf{Sortie} & \textbf{Fautes métier} \\ \midrule
\endfirsthead
\toprule \textbf{Opération} & \textbf{Entrée} & \textbf{Sortie} & \textbf{Fautes métier} \\ \midrule
\endhead
\bottomrule
\endlastfoot
\textbf{\texttt{creerOffre}}\newline {\footnotesize Crée une offre après avoir vérifié l'entreprise auprès de S2.} & \texttt{offre} : Offre & \texttt{offreCreee} : Offre & \texttt{DonneesInvalides} \newline \texttt{ServiceIndisponible} \\ \midrule
\textbf{\texttt{modifierOffre}}\newline {\footnotesize Remplace l'offre. Le statut suit les transitions autorisées ; l'entreprise ne change pas.} & \texttt{id} : long \newline \texttt{offre} : Offre & \texttt{offreModifiee} : Offre & \texttt{OffreIntrouvable} \newline \texttt{DonneesInvalides} \newline \texttt{Conflit} \\ \midrule
\textbf{\texttt{supprimerOffre}}\newline {\footnotesize Supprime une offre en BROUILLON. Une offre publiée se clôture.} & \texttt{id} : long & – & \texttt{OffreIntrouvable} \newline \texttt{Conflit} \\ \midrule
\textbf{\texttt{consulterOffre}}\newline {\footnotesize Renvoie une offre.} & \texttt{id} : long & \texttt{offre} : Offre & \texttt{OffreIntrouvable} \\ \midrule
\textbf{\texttt{listerOffres}}\newline {\footnotesize Liste les offres selon des critères facultatifs.} & \texttt{filtre} : FiltreOffres, facultatif & \texttt{offre} : Offre[], facultatif & \texttt{DonneesInvalides} \\ \midrule
\textbf{\texttt{verifierDisponibilite}}\newline {\footnotesize Vérifie que l'offre est active (R2), avant sa date limite (R3) et avec une place (R4). Une offre inconnue donne existe=false, pas une faute.} & \texttt{offreId} : long & \texttt{disponibilite} : Disponibilite & \texttt{ServiceIndisponible} \\ \midrule
\textbf{\texttt{verifierCompatibilite}}\newline {\footnotesize Compare l'offre au profil de l'étudiant obtenu auprès de S1 (R7, R3) et calcule le score indicatif.} & \texttt{offreId} : long \newline \texttt{etudiantId} : long & \texttt{compatibilite} : Compatibilite & \texttt{OffreIntrouvable} \newline \texttt{EtudiantIntrouvable} \newline \texttt{ServiceIndisponible} \\ \midrule
\textbf{\texttt{decrementerPlaces}}\newline {\footnotesize Retire une place de façon atomique. L'offre passe en POURVUE à zéro (R10).} & \texttt{offreId} : long & \texttt{offre} : Offre & \texttt{OffreIntrouvable} \newline \texttt{Conflit} \\
\end{longtable}
}

\section{Extrait de WSDL}

Le listing~\ref{lst:wsdl} montre la définition de \code{verifierEtudiant} : le type de la réponse, puis l'opération dans le \code{portType} avec ses messages et ses fautes.

\lstinputlisting[language=XMLfr, caption={Extrait de \code{etudiants.wsdl} : opération \code{verifierEtudiant}}, label={lst:wsdl}]{generes/extrait-wsdl-etudiants.xml}

\section{Exemples d'échanges}

Les échanges suivants ont été enregistrés sur les services en fonctionnement par \code{rapport/outils/capturer\_exemples.py}.

\subsection{Vérification d'un étudiant suspendu}
\lstinputlisting[language=XMLfr, caption={Requête \code{verifierEtudiant} avec en-tête WS-Security et corrélation}]{generes/exemples/soap-verifier-suspendu-requete.xml}
\lstinputlisting[language=XMLfr, caption={Réponse : l'étudiante existe mais n'est pas éligible (règle R1)}]{generes/exemples/soap-verifier-suspendu-reponse.xml}

\subsection{Faute typée}
Pour un identifiant inconnu, \code{consulterEtudiant} renvoie la faute déclarée dans le WSDL, avec son détail structuré. Conformément à SOAP~1.1, le statut HTTP d'une faute est 500.
\lstinputlisting[language=XMLfr, caption={Réponse à \code{consulterEtudiant(9999)} : \code{EtudiantIntrouvableFault}}]{generes/exemples/soap-consulter-inconnu-reponse.xml}

Sans en-tête WS-Security, la même requête est refusée avant d'atteindre le service, par une faute \code{wsse:InvalidSecurity} (tableau~\ref{tab:wsse}).

\section{Implémentation}

Chaque service SOAP comprend quatre couches : l'\emph{endpoint} JAX-WS généré à partir du WSDL, qui ne fait que convertir les types ; un service métier qui applique les règles et les transitions de statut ; un DAO JDBC ; et des clients SOAP vers les autres services quand c'est nécessaire (S2 vers S3, S3 vers S1 et S2). Les gestionnaires JAX-WS (\emph{handlers}) communs, définis dans le module \code{common}, vérifient le jeton WS-Security, reprennent l'identifiant de corrélation et journalisent chaque appel avec le compte qui l'a fait.

Le serveur utilisé est celui intégré à la JDK (\code{com.sun.net.httpserver}) piloté par Metro : aucun serveur d'applications n'est nécessaire, chaque service démarre en une à deux secondes.
